% Figure: Campbell diagram for a rotating shaft. Horizontal axis is
% rotation speed; vertical axis is frequency. Horizontal lines are the
% (roughly speed-independent) shaft bending natural frequencies; the
% diagonal lines are engine-order excitation lines (1x, 2x). The
% intersections mark critical speeds where an excitation order meets a
% natural frequency.
\begin{figure}[htbp]
\centering
\begin{tikzpicture}
\begin{axis}[
  width=12cm, height=7.2cm,
  xlabel={rotation speed $N$ (rpm)},
  ylabel={frequency (Hz)},
  xmin=0, xmax=6000, ymin=0, ymax=180,
  axis lines=left,
  every axis x label/.style={at={(current axis.right of origin)}, anchor=west},
  every axis y label/.style={at={(current axis.above origin)}, anchor=south},
  grid=major, grid style={enggray!20},
  tick label style={font=\scriptsize},
  label style={font=\small},
  legend style={font=\scriptsize},
  legend pos=north west,
]
% first two bending natural frequencies (horizontal), drawn as plots
\addplot[engblue, very thick, domain=0:6000, samples=2] {42};
\addlegendentry{1st bending mode}
\addplot[engblue, very thick, dashed, domain=0:6000, samples=2] {120};
\addlegendentry{2nd bending mode}
% engine-order lines: f = order * N / 60
\addplot[engred, thick, domain=0:6000, samples=2] {x/60};
\addlegendentry{$1\times$ order}
\addplot[engorange, thick, domain=0:6000, samples=2] {2*x/60};
\addlegendentry{$2\times$ order}
% critical-speed intersections (1x meets 1st mode at N=2520; 2x meets 1st mode at N=1260)
\addplot[engblack, only marks, mark=*, mark size=2pt, forget plot] coordinates {(2520,42) (1260,42)};
\node[engblack, font=\scriptsize, anchor=south west] at (axis cs:2520,44) {1st critical};
\node[engblack, font=\scriptsize, anchor=south east] at (axis cs:1260,44) {$2\times$ crossing};
% operating band shading via two dotted verticals
\addplot[enggray, dotted, thick, forget plot] coordinates {(3000,0) (3000,180)};
\addplot[enggray, dotted, thick, forget plot] coordinates {(3600,0) (3600,180)};
\node[enggray, font=\scriptsize, anchor=south] at (axis cs:3300,160) {operating band};
\end{axis}
\end{tikzpicture}
\caption[Campbell diagram of a rotating shaft]{Campbell diagram for a
rotating shaft, the standard tool for placing critical speeds. The
horizontal lines are the shaft bending natural frequencies, nearly
independent of speed for a stiff rotor; the diagonal lines are the
excitation orders, whose frequency rises in proportion to rotation
speed ($1\times$ at the running speed, $2\times$ at twice it). Each
intersection of an order with a natural frequency marks a speed at
which that order resonates with that mode, a critical speed. The first
critical here, where the $1\times$ order meets the first bending mode,
falls at $2520\,\text{rpm}$, safely below the operating band of $3000$
to $3600\,\text{rpm}$ marked by the dotted lines, so the shaft runs
supercritically with respect to its first mode while the second mode
stays well above any order in the band. Design discipline is to read
this diagram and keep the operating band clear of every intersection
by a margin.}
\label{fig:vol03ch06:campbell}
\end{figure}
